% Minimal retained preamble for lane-p2 cards taken from
% sources/logic-fl02-conservation-of-ramsey-s-theorem-f/pinned-original-source.tex
% (Conservation of Ramsey's theorem for pairs, arXiv:2402.11616v3).
%
% Copied verbatim from the author's own file:
%   line 2      : \documentclass{article}
%   lines 16-25 : the author's theorem declarations
%   lines 46-88 : the author's macro definitions for the reverse-mathematics principles
%   line 91     : \card
%   lines 104-107: the natural sums and products, which rely on mathabx (source line 103)
%   line 116    : \qt
%
% Omitted author preamble lines and why (nothing was renamed or redefined):
%   line 3 (graphicx), 7 (url), 9 (cancel): no figure, url or cancel use in the retained spans;
%   line 6 (cleveref): the assembler adds its own cleveref machinery, a second load clashes;
%   line 8 (\usepackage[utf8x]{inputenc}): inputenc with the utf8x option fights the forced
%       xelatex engine;
%   lines 28-37 (the rep@theorem machinery and \newreptheorem{maintheorem}): the retained spans
%       use no rep-theorem environment;
%   lines 40-43 (xcolor and the three author comment macros): not used by the retained spans;
%   lines 93-99 (the LS2 font encoding block for \smashtimes based on stix): the host TeX Live has
%       no stix directory and the author himself never uses the result -- line 105 maps \ntimes to
%       \dottimes, and line 110 shows \smashtimes already commented out in the source.
% mathabx is retained because \dotplus, \dottimes and \dotdiv come from it.
\documentclass{article}
\usepackage{amsmath,amssymb,amsthm}
\usepackage{mathabx}
\usepackage{cancel}
\usepackage{cleveref}

\newtheorem{theorem}{Theorem}
\numberwithin{theorem}{section}
\newtheorem{maintheorem}[theorem]{Main Theorem}
\newtheorem{lemma}[theorem]{Lemma}
\newtheorem{question}[theorem]{Question}
\newtheorem{proposition}[theorem]{Proposition}
\newtheorem{statement}[theorem]{Statement}
\newtheorem{remark}[theorem]{Remark}
\newtheorem{definition}[theorem]{Definition}
\newtheorem{corollary}[theorem]{Corollary}

\newcommand{\RCA}{\mathsf{RCA}}
\newcommand{\WKL}{\mathsf{WKL}}
\newcommand{\SRT}{\mathsf{SRT}}
\newcommand{\EM}{\mathsf{EM}}
\newcommand{\SEM}{\mathsf{SEM}}
\newcommand{\ADS}{\mathsf{ADS}}
\newcommand{\SCAC}{\mathsf{SCAC}}
\newcommand{\SADS}{\mathsf{SADS}}
\newcommand{\COH}{\mathsf{COH}}
\newcommand{\ACA}{\mathsf{ACA}}
\newcommand{\qvdash}{\operatorname{{?}{\vdash}}}
\newcommand{\nqvdash}{\operatorname{{?}{\nvdash}}}
\newcommand{\BSig}{\mathsf{B}\Sigma^0}
\newcommand{\ISig}{\mathsf{I}\Sigma^0}
\newcommand{\LSig}{\mathsf{L}\Sigma^0}
\newcommand{\BPi}{\mathsf{B}\Pi^0}
\newcommand{\IPi}{\mathsf{I}\Pi^0}
\newcommand{\LPi}{\mathsf{L}\Pi^0}
\newcommand{\IDel}{\mathsf{I}\Delta^0}
\newcommand{\BCDe}{\mathsf{BC}\Delta^0}
\newcommand{\WF}{\mathsf{WF}}
\newcommand{\RT}{\mathsf{RT}}
\newcommand{\BME}{\mathsf{BME}}
\newcommand{\TT}{\mathsf{TT}}
\renewcommand{\O}{\mathcal{O}}
\renewcommand{\L}{\mathcal{L}}
\newcommand{\C}{\mathcal{C}}
\newcommand{\R}{\mathcal{R}}
\renewcommand{\S}{\mathcal{S}}
\newcommand{\T}{\mathcal{T}}
\renewcommand{\P}{\mathcal{P}}
\newcommand{\M}{\mathcal{M}}
\newcommand{\Nc}{\mathcal{N}}
\renewcommand{\P}{\mathcal{P}}
\newcommand{\Q}{\mathcal{Q}}
\newcommand{\NN}{\mathbb{N}}
\newcommand{\PP}{\mathbb{P}}
\newcommand{\uh}[0]{{\upharpoonright}}
\newcommand{\card}{\operatorname{card}}
\newcommand{\nplus}{\dotplus}
\newcommand{\ntimes}{\dottimes}
\newcommand{\nminus}{\dotdiv}
\newcommand{\nsum}{\dot{\sum}}
\def\qt#1{``#1''}
\title{Conservation transfers from omega\_k powers to the union over k of omega to the omega\_k}
\author{Quentin Le Houérou, Ludovic Levy Patey, Keita Yokoyama}
\date{arxiv:2402.11616v3; CC BY 4.0}
\begin{document}
\maketitle
\noindent\textit{Editorial note.} The counted unit is the article's own Theorem theorem:problem-wf-cons-2 at source lines 1425-1428 together with the proof environment that begins two lines after the statement ends (1430-1440); the pairing is the ordinary adjacent one.\par\medskip
\noindent\textit{Editorial note.} The retained context is only the subsection heading at source line 1421, which names the weakening of the well-foundedness assumption; the explanatory paragraph that follows it at source line 1423 is not retained because it carries a citation of Chong, Slaman and Yang, and the result it describes is not needed to read the counted theorem.\par\medskip
\noindent\textit{Editorial note.} Two neighbouring items are deliberately not counted: the repeated main theorem at source lines 1411-1414 and the one at 1442-1445 use the author's custom repeated-theorem environment and their proofs at 1416-1418 and 1447-1449 are two-line chains of applications without independent content, which falls under the follows-easily exclusion; the commented-out proposition at source lines 1404-1409 is not author text for our purposes.\par\medskip
\noindent\textit{Editorial note.} The retained preamble reproduces verbatim the author's document class at source line 2, his theorem declarations at source lines 16-25, his macro definitions at source lines 46-88 and 91, and his natural-operation macros at source lines 104-107, and it retains mathabx and cleveref; the stix-dependent LS2 block and the inputenc load with the utf8x option are not retained, and the repeated-theorem machinery at source lines 28-37 is not retained because no retained span uses it.\par\medskip
\noindent\textit{Editorial note.} No retained line contains a citation, so the card carries no bibliography.\par\medskip
\noindent\textit{Allowed standard prerequisites (editorial).}\begin{itemize}
\item Theories RCA\_0, the schemes T\_n and the well-foundedness statements. RCA\_0 is the base theory of reverse mathematics, T\_n is either collection or induction for Sigma\textasciicircum{}0\_n formulas, and WF of an ordinal states that the ordinal is well founded; the notation WF(omega\_k\textasciicircum{}alpha) refers to the ordinal omega with subscript k and exponent alpha, and the union over standard k of WF(omega to the omega\_k) is the axiom system considered in this section; the article uses these notations throughout, for instance in its introduction.. They are the theories appearing in the hypothesis and the conclusion of the counted theorem, and in the assumption of its proof.
\item Compactness and omega-extensions of topped models. The proof extends the language by a constant, uses that a theory is finitely satisfiable to obtain a model by compactness, and iterates omega-extensions of countable topped models; these are standard tools of the article, used the same way in the earlier conservation theorem of the paper.. The proof obtains a model with a non-standard constant by compactness and then builds an increasing sequence of omega-extensions whose union satisfies the theory together with a counterexample.
\item Pi\textasciicircum{}1\_1 conservativity. A problem is Pi\textasciicircum{}1\_1 conservative over a theory when every arithmetical sentence provable in the theory together with the problem is already provable in the theory; the proof argues by contraposition on arithmetical sentences.. The counted theorem asserts conservativity of a Pi\textasciicircum{}1\_2 problem, and the proof establishes it by refuting provability of an arithmetical sentence.
\end{itemize}\medskip
\subsection{Preserving $\bigcup_n \WF(\omega_n)$}\label[section]{sect:bigcup-wf-omegan}

\begin{theorem}\label[theorem]{theorem:problem-wf-cons-2}
Fix~$n \geq 2$. Let~$T_n$ be either~$\BSig_n$ or~$\ISig_n$.
Let $\mathsf{P}$ be a $\Pi_2^1$ problem such that for every countable topped model $\M = (M,S) \models \RCA_0 + T_n + \WF(\omega_k^{\alpha})$ with $\alpha \in M$ an ordinal less than $\epsilon_0$ and $k$ a standard integer and every instance of $\mathsf{P}$ in $\M$ there exists a solution $G \subset M$ of that instance such that $\M[G] \models \RCA_0 + T_n + \WF(\alpha)$. Then $\mathsf{P}$ is $\Pi_1^1$-conservative over $\RCA_0 + T_n + \bigcup_{k \in \omega} \WF(\omega^\omega_k)$.   
\end{theorem}

\begin{proof}
Assume $\RCA_0 + T_n + \bigcup_{n \in \omega} \WF(\omega_n) \not \vdash \forall X \phi(X)$ for $\phi$ arithmetical.  

Extend the language of second order arithmetics with a first-order constant symbol $c$. The theory $\RCA_0 + T_n + \WF(\omega_c) + \bigcup_{k \in \omega} c > k + \exists X \neg \phi(X)$ is finitely satisfiable by our assumption, so by compactness we have a model $\M \models \RCA_0 + T_n + \WF(\omega_c) + \bigcup_{k \in \omega} c > k + \exists X \neg \phi(X)$. Such a $c$ will therefore be non-standard.

From our assumption, for any countable topped model $\mathcal{N} = (N,S) \models \RCA_0 + T_n + \WF(\omega_{d+k})$ and instance of $\mathsf{P}$ in this model, we can find a countable topped $\omega$-extension that will be a model of $\RCA_0 + T_n + \WF(\omega_{d})$ and will contain a solution to that instance.

So by iterating that reasoning, we can build a sequence  $\M_0 = \M \subseteq \M_1 \subseteq \dots$ of countable topped model such that $\M_s \models \RCA_0 + T_n + \WF(\omega_{c - ks})$ and such that every instance of $P$ appearing inside one $\M_i$ will eventually have a solution in one $\M_j$. Since $c$ is non-standard and every such $s$ is standard, every $\M_i$ is a model of $\bigcup_{k \in \omega} \WF(\omega^\omega_k)$, and the model $\bigcup_{s \in \omega} \M_s$ will therefore be a model of $\RCA_0 + T_n + \bigcup_{k \in \omega} \WF(\omega^\omega_k) + \exists X \neg \phi(X) + \mathsf{P}$.

So $\RCA_0 + T_n + \bigcup_{k \in \omega} \WF(\omega^\omega_k)  + \mathsf{P} \not \vdash \forall X \phi(X)$.
\end{proof}

\end{document}
